% =========================================================
% Proof Stub: Equivalence of Convergence Definitions
% Source: volume-iii/analysis/sequences/notes/sequences/notes-convergence.tex
% Mode: standard
% =========================================================

\clearpage
\phantomsection
\hypertarget{proof-RA-ABB-T-equiv-conv-defns}{}

\subsubsection[Equivalence of Convergence Definitions]{Proof --- Equivalence of Convergence Definitions}

\bigskip

\noindent
\textbf{Source.}
See \texttt{notes-convergence.tex}.

\vspace{0.75em}

\noindent
\textbf{Goal.}
In $\mathbb{R}$, the $\varepsilon$-definition of convergence and the neighborhood definition are equivalent.

\vspace{0.75em}

\noindent
\textbf{Logical form.}
\[
  (x_n) \to L \text{ in } \varepsilon\text{-sense} \Longleftrightarrow \forall \text{ open } U \ni L,\; \exists N,\; \forall n \ge N,\; x_n \in U.
\]

\vspace{0.75em}

\noindent
\textbf{Proof strategy.}
For the forward direction: given open $U \ni L$, find $\varepsilon > 0$ with $(L-\varepsilon, L+\varepsilon) \subseteq U$, then use the $\varepsilon$-$N$ condition. For the backward direction: $(L-\varepsilon, L+\varepsilon)$ is open, apply the hypothesis.

\vspace{0.75em}

\noindent
\textbf{Proof.}
\begin{proof}
\Claim In $\mathbb{R}$, the $\varepsilon$-definition of convergence and the neighborhood definition are equivalent.

\Given 

\Goal 

\Strategy For the forward direction: given open $U \ni L$, find $\varepsilon > 0$ with $(L-\varepsilon, L+\varepsilon) \subseteq U$, then use the $\varepsilon$-$N$ condition. For the backward direction: $(L-\varepsilon, L+\varepsilon)$ is open, apply the hypothesis.

\medskip

% --- (⇒) given open U ∋ L, find ε with (L-ε, L+ε) ⊆ U (open set in R contains a ball) ---
% --- apply ε-N definition to get N; for n ≥ N: x_n ∈ (L-ε,L+ε) ⊆ U ---
% --- (⇐) given ε > 0, the set (L-ε,L+ε) is open and contains L ---
% --- apply neighborhood hypothesis to U = (L-ε,L+ε); get N such that x_n ∈ U for n ≥ N ---

\end{proof}

\begin{remark*}[Proof shape]
Biconditional: each direction follows by choosing the right $\varepsilon$ or neighborhood.
\end{remark*}

\begin{remark*}[Dependencies]
\begin{itemize}
  \item Open sets in $\mathbb{R}$ contain $\varepsilon$-balls around each of their points.
  \item $\varepsilon$-$N$ definition of convergence.
  \item $(L-\varepsilon, L+\varepsilon)$ is an open neighborhood of $L$ for every $\varepsilon > 0$.
\end{itemize}
\end{remark*}
